%!TEX program = lualatex
\documentclass[12pt]{article}                           % font size and type of document
\usepackage[a4paper, margin=25mm]{geometry}             % paper size and margins
\usepackage[onehalfspacing]{setspace}                   % line spacing
\usepackage{lipsum}                                     % create dummy text in order to check
                                                        % formating

\usepackage[T1]{fontenc}                                % hyphenating of words containing accented   
                                                                                                        % chars, font encoding for accented chars
\usepackage{textcomp}                                   % additional symbols
%\usepackage{mathptmx}                                   % font type (Times)

\usepackage{fontspec}
\setmainfont{Trebuchet MS}

%\usepackage{lmodern}                                    % font type (Latin Modern Sans Serif)
%\usepackage{tgadventor}                                 % font type (TEX Gyre Adventor)
\usepackage[utf8]{inputenc}                             % used for mutated vowels e.g. ä, ü, ö, ß
\usepackage[ngerman]{babel}                             % load multilingual support for german
\usepackage[babel,german=guillemets]{csquotes}          % ensure multilingual support for biblatex 
                                                        % quoting
\newcommand{\titlename}{Projektierungshandbuch TIA-Portal}                    % change title of document
\newcommand{\version}{15.0}  
\newcommand{\authorname}{Daniel Klein-Günnewick}               % and name of author here
\newcommand{\state}{Verbindlich, in Bearbeitung}
\newcommand{\lasteditor}{Stephan Kampshoff}
\newcommand{\lebbing}{LEBBING AUTOMATION \& DRIVES GMBH}

                                                        % applied in whole document
\title{\titlename}
\author{\authorname}
\date{\today}
\newcounter{savepage}

\usepackage[hyphens]{url}                               % hyphening links
\usepackage{hyperref}
\hypersetup{
    colorlinks=true,
    linkcolor=black,
    citecolor=black,
    filecolor=magenta,      
    urlcolor=cyan,
    pdftitle={\titlename},
    pdfauthor={\authorname}
    }
\urlstyle{same}                                         % allows links within pdf document
                                                        % and to external websites

\usepackage{xcolor}                                     % load xcolor package for customization
\usepackage{listings}                                   % supports listings environments
\definecolor{codegreen}{rgb}{0,0.6,0}                   % define and set custom colors
\definecolor{codegray}{rgb}{0.5,0.5,0.5}
\definecolor{codepurple}{rgb}{0.58,0,0.82}
\definecolor{backcolour}{rgb}{0.95,0.95,0.92}

\lstdefinestyle{mystyle}{                               % define listings style
    language=[LaTeX]TeX,
    morecomment=[l]{//},
    texcsstyle=*[1]\color{codepurple},
    emphstyle=\color{codepurple},
    backgroundcolor=\color{backcolour},   
    commentstyle=\color{black},   
    basicstyle=\ttfamily\footnotesize,
    breakatwhitespace=false,         
    breaklines=true,                 
    captionpos=b,                    
    keepspaces=true,                 
    numbers=left,                    
    numbersep=5pt,                  
    showspaces=false,                
    showstringspaces=false,
    showtabs=false,                  
    tabsize=2
}


\renewcommand{\lstlistlistingname}{Quellcodeverzeichnis}
\renewcommand{\lstlistingname}{Quellcode}               % change name of list of listings
\lstset{style=mystyle}                                  % set listings style to my defined style

\usepackage{graphicx}                                   % required for inserting images
\graphicspath{{images/}}
\usepackage{caption}                                    % more customization of captions
%\captionsetup{figurename=Abb., tablename=Tab.}
\usepackage{wrapfig}                                    % wrap text around figures
%\DeclareCaptionType{mycapequ}[][Formelverzeichnis]      % adds list of equations
%\DeclareCaptionLabelFormat{nan}{Gleichung #2}           % caption compatibility with equations 
%\captionsetup[mycapequ]{labelformat=nan}

\usepackage[final]{pdfpages}                            % adds support for pdf file integration
\usepackage{tikz}                                       % create graphics in latex environment
\usetikzlibrary{positioning, calc}                      % enable relative positioning of nodes in tikz figure
\usepackage{pgfplots}                                   % create diagram, charts etc. with tikz library
\pgfplotsset{compat=1.18}                               % set used pgfplot package version for compatibility
\usepackage{amsmath}                                    % enhanced mathematical expressions
\usepackage{siunitx}                                    % includes si units support
\usepackage[ngerman, noabbrev]{cleveref}                % improves referencing of tables, 
                                                        % figures, equations, etc.
%\crefformat{equation}{Gl.~(#2#1#3)}                    % customise reference to Gl. (default
                                                        %  "Gleichung")
\crefname{listing}{Quellcode}{Quellcodes}               % change reference name
\Crefname{listing}{Quellcode}{Quellcodes}

\usepackage{tocloft}                                    % extension for toc, lof, lot and other lists
\renewcommand{\cftsecleader}{\cftdotfill{\cftdotsep}}   % dot line for sections in toc
\usepackage{tocbibind}                                  % adds all lists to table of contents

\usepackage[acronym,toc,translate=babel,
            nonumberlist, nopostdot]{glossaries}
\renewcommand{\glsnamefont}[1]{\textbf{#1}}
\setlength\LTleft{0pt}
\setlength\LTright{0pt}
\setlength\glsdescwidth{0.8\hsize}                      % adds support for glossary and
%\makeglossaries                                        % acronyms and abbreviations

\usepackage[                                            % header for bibliography
    backend=biber, 
    natbib=true,
    hyperref=true,
    style=ieee
]{biblatex}
%\DeclareLanguageMapping{ngerman}{german-apa}           % only necessary for biblatex-apa style
%\DeclareFieldInputHandler{extradate}{\def\NewValue{}}  % remove suffix letters if author and year
\addbibresource{references.bib}                         % are identical in bibliography

%Abbildungsverzeichnis
\usepackage{tocloft}

\usepackage{tabto}

\usepackage{fancyhdr}   %Kopfzeile
\makeatletter

% Kopfzeilenstil definieren
\pagestyle{fancy}
\fancyhf{} % Alle Kopf- und Fußzeilen löschen
\renewcommand{\headrulewidth}{0pt} % Horizontale Linie in Kopfzeile entfernen
\renewcommand{\sectionmark}[1]{\markright{\thesection\ #1}}
% Kopfzeile definieren
\fancyhead[L]{\includegraphics[width=3.8cm]{LebbingLogo.png}} % Bild in der linken Kopfzeile
\fancyhead[R]{ \titlename~ V\version \par Status: \state \par \rightmark} % Rechter Kopfzeilentext
% Höhe der Kopfzeile festlegen
\setlength{\headheight}{30pt} % Größerer Wert für mehr Platz

%Fußzeile definieren
\fancyfoot[L]{\lebbing}
\fancyfoot[R]{\thepage}

\makeatother

\usepackage{tabulary} %Einstellungen von Tabellen
\usepackage{array}
\setlength{\tabcolsep}{10pt}            %Einstellen des abstands nach links und rechts in einer zeile
\renewcommand{\arraystretch}{1.35}       %Einstellen der Zeilenhöhe
\usepackage{enumitem} % Zum Anpassen der Aufzählungseigenschaften

\usepackage{makecell}
\usepackage{afterpage}                                  % commands get expanded after the current                                                           % page  
\newcommand\blankpage{                                  % is output using \afterpage{\command}
    \newpage                                            % -> useful for \clearpage
    \null                                               % define command to insert blank page
    \thispagestyle{empty}
    \newpage}
